% Options for packages loaded elsewhere
\PassOptionsToPackage{unicode}{hyperref}
\PassOptionsToPackage{hyphens}{url}
%


\PassOptionsToPackage{table}{xcolor}

\documentclass[
  10pt,
  letterpaper,
]{article}

\usepackage{amsmath,amssymb}
\usepackage{iftex}
\ifPDFTeX
  \usepackage[T1]{fontenc}
  \usepackage[utf8]{inputenc}
  \usepackage{textcomp} % provide euro and other symbols
\else % if luatex or xetex
  \usepackage{unicode-math}
  \defaultfontfeatures{Scale=MatchLowercase}
  \defaultfontfeatures[\rmfamily]{Ligatures=TeX,Scale=1}
\fi
\usepackage{lmodern}
\ifPDFTeX\else  
    % xetex/luatex font selection
\fi
% Use upquote if available, for straight quotes in verbatim environments
\IfFileExists{upquote.sty}{\usepackage{upquote}}{}
\IfFileExists{microtype.sty}{% use microtype if available
  \usepackage[]{microtype}
  \UseMicrotypeSet[protrusion]{basicmath} % disable protrusion for tt fonts
}{}
\makeatletter
\@ifundefined{KOMAClassName}{% if non-KOMA class
  \IfFileExists{parskip.sty}{%
    \usepackage{parskip}
  }{% else
    \setlength{\parindent}{0pt}
    \setlength{\parskip}{6pt plus 2pt minus 1pt}}
}{% if KOMA class
  \KOMAoptions{parskip=half}}
\makeatother
\usepackage{xcolor}
\usepackage[top=0.85in,left=2.75in,footskip=0.75in]{geometry}
\setlength{\emergencystretch}{3em} % prevent overfull lines
\setcounter{secnumdepth}{-\maxdimen} % remove section numbering


\providecommand{\tightlist}{%
  \setlength{\itemsep}{0pt}\setlength{\parskip}{0pt}}\usepackage{longtable,booktabs,array}
\usepackage{calc} % for calculating minipage widths
% Correct order of tables after \paragraph or \subparagraph
\usepackage{etoolbox}
\makeatletter
\patchcmd\longtable{\par}{\if@noskipsec\mbox{}\fi\par}{}{}
\makeatother
% Allow footnotes in longtable head/foot
\IfFileExists{footnotehyper.sty}{\usepackage{footnotehyper}}{\usepackage{footnote}}
\makesavenoteenv{longtable}
\usepackage{graphicx}
\makeatletter
\newsavebox\pandoc@box
\newcommand*\pandocbounded[1]{% scales image to fit in text height/width
  \sbox\pandoc@box{#1}%
  \Gscale@div\@tempa{\textheight}{\dimexpr\ht\pandoc@box+\dp\pandoc@box\relax}%
  \Gscale@div\@tempb{\linewidth}{\wd\pandoc@box}%
  \ifdim\@tempb\p@<\@tempa\p@\let\@tempa\@tempb\fi% select the smaller of both
  \ifdim\@tempa\p@<\p@\scalebox{\@tempa}{\usebox\pandoc@box}%
  \else\usebox{\pandoc@box}%
  \fi%
}
% Set default figure placement to htbp
\def\fps@figure{htbp}
\makeatother

% Use adjustwidth environment to exceed column width (see example table in text)
\usepackage{changepage}

% marvosym package for additional characters
\usepackage{marvosym}

% cite package, to clean up citations in the main text. Do not remove.
% Using natbib instead
% \usepackage{cite}

% Use nameref to cite supporting information files (see Supporting Information section for more info)
\usepackage{nameref,hyperref}

% line numbers
\usepackage[right]{lineno}

% ligatures disabled
\usepackage{microtype}
\DisableLigatures[f]{encoding = *, family = * }

% create "+" rule type for thick vertical lines
\newcolumntype{+}{!{\vrule width 2pt}}

% create \thickcline for thick horizontal lines of variable length
\newlength\savedwidth
\newcommand\thickcline[1]{%
  \noalign{\global\savedwidth\arrayrulewidth\global\arrayrulewidth 2pt}%
  \cline{#1}%
  \noalign{\vskip\arrayrulewidth}%
  \noalign{\global\arrayrulewidth\savedwidth}%
}

% \thickhline command for thick horizontal lines that span the table
\newcommand\thickhline{\noalign{\global\savedwidth\arrayrulewidth\global\arrayrulewidth 2pt}%
\hline
\noalign{\global\arrayrulewidth\savedwidth}}

% Text layout
\raggedright
\setlength{\parindent}{0.5cm}
\textwidth 5.25in 
\textheight 8.75in

% Bold the 'Figure #' in the caption and separate it from the title/caption with a period
% Captions will be left justified
\usepackage[aboveskip=1pt,labelfont=bf,labelsep=period,justification=raggedright,singlelinecheck=off]{caption}
\renewcommand{\figurename}{Fig}

% Remove brackets from numbering in List of References
\makeatletter
\renewcommand{\@biblabel}[1]{\quad#1.}
\makeatother

% Header and Footer with logo
\usepackage{lastpage,fancyhdr}
\usepackage{epstopdf}
%\pagestyle{myheadings}
\pagestyle{fancy}
\fancyhf{}
%\setlength{\headheight}{27.023pt}
%\lhead{\includegraphics[width=2.0in]{PLOS-submission.eps}}
\rfoot{\thepage/\pageref{LastPage}}
\renewcommand{\headrulewidth}{0pt}
\renewcommand{\footrule}{\hrule height 2pt \vspace{2mm}}
\fancyheadoffset[L]{2.25in}
\fancyfootoffset[L]{2.25in}
\lfoot{\today}
\usepackage{booktabs}
\usepackage{caption}
\usepackage{longtable}
\usepackage{colortbl}
\usepackage{array}
\usepackage{anyfontsize}
\usepackage{multirow}
\usepackage{setspace}
\doublespacing
\makeatletter
\@ifpackageloaded{caption}{}{\usepackage{caption}}
\AtBeginDocument{%
\ifdefined\contentsname
  \renewcommand*\contentsname{Table of contents}
\else
  \newcommand\contentsname{Table of contents}
\fi
\ifdefined\listfigurename
  \renewcommand*\listfigurename{List of Figures}
\else
  \newcommand\listfigurename{List of Figures}
\fi
\ifdefined\listtablename
  \renewcommand*\listtablename{List of Tables}
\else
  \newcommand\listtablename{List of Tables}
\fi
\ifdefined\figurename
  \renewcommand*\figurename{Fig}
\else
  \newcommand\figurename{Fig}
\fi
\ifdefined\tablename
  \renewcommand*\tablename{Table}
\else
  \newcommand\tablename{Table}
\fi
}
\@ifpackageloaded{float}{}{\usepackage{float}}
\floatstyle{ruled}
\@ifundefined{c@chapter}{\newfloat{codelisting}{h}{lop}}{\newfloat{codelisting}{h}{lop}[chapter]}
\floatname{codelisting}{Listing}
\newcommand*\listoflistings{\listof{codelisting}{List of Listings}}
\makeatother
\makeatletter
\makeatother
\makeatletter
\@ifpackageloaded{caption}{}{\usepackage{caption}}
\@ifpackageloaded{subcaption}{}{\usepackage{subcaption}}
\makeatother

\usepackage[numbers,square,comma]{natbib}
\bibliographystyle{plos2025}
\usepackage{bookmark}

\IfFileExists{xurl.sty}{\usepackage{xurl}}{} % add URL line breaks if available
\urlstyle{same} % disable monospaced font for URLs
\hypersetup{
  pdftitle={Periodicity in Steller's eider (Polysticta stelleri) population size and density on the Arctic Coastal Plain, Alaska, revealed using generalized additive models},
  pdfauthor={Erik E. Osnas},
  hidelinks,
  pdfcreator={LaTeX via pandoc}}

\begin{document}
\vspace*{0.2in}

% Title must be 250 characters or less.
\begin{flushleft}
{\Large
\textbf\newline{Periodicity in Steller's eider (\emph{Polysticta
stelleri}) population size and density on the Arctic Coastal Plain,
Alaska, revealed using generalized additive
models} % Please use "sentence case" for title and headings (capitalize only the first word in a title (or heading), the first word in a subtitle (or subheading), and any proper nouns).
}
\newline
\\
% Insert author names, affiliations and corresponding author email (do not include titles, positions, or degrees).
Erik E. Osnas\textsuperscript{1*}
\\
\bigskip
\textsuperscript{1} Division of Migratory Bird Management, United States Fish and Wildlife Service, Anchorage, Alaska, United States of America 
\bigskip

* Corresponding author \newline
E-mail: erik\_osnas@fws.gov (EO)

\end{flushleft}

\section*{Abstract}
Annual population and spatial density estimates are needed for the
threatened Alaska-breeding population of Steller's eiders
(\emph{Polysticta stelleri}) to assess recovery status and guide
recovery actions but these quantities are poorly estimated in this rare
species using traditional methods. Therefore, population size and
spatial variation in density were estimated across the Arctic Coastal
Plain (ACP) and Utqiagvik Triangle (Triangle) survey areas using
spatio-temporal generalized additive models and compared to traditional
design-based estimates. Compared to design-based estimates, model-based
estimates were more precise and less variable between years. Moreover,
data sets can be combined to produce common estimates, detection
probability can be incorporated to estimate population size, and spatial
density maps can be produced from model-based estimates. Across both
data sets and when combined, Steller's eider populations fluctuated with
an approximate 6.53-year period (95\% CI: 5.46, 7.93) with populations
cycling from modelled posterior lows of 25 to 500 individuals and highs
from 230 to \textgreater 2000 individuals across the ACP. Over the long
term, the posterior 25-year geometric mean growth rate was -0.02 (95\% CI:
-0.07, 0.02), and shorter term growth rates were less well-estimated and
fluctuated between positive and near zero. Density maps showed a high
concentration in the northern part of the Triangle area and lower
densities southward and across the ACP. Models fit to just the Triangle
area indicated that eider density has been shifting northward though
time, but this pattern was not supported in the sparse ACP data set, nor
was it well-estimated when the ACP and Triangle data sets were combined.
Given the strong cycle in this population, estimates of population size
at similar cycle phases should be used for trend estimates.

\linenumbers

\section{Introduction}\label{introduction}

The Pacific population of Steller's eider (\emph{Polysticta stelleri})
is a seaduck that nests in northeastern Russia and Alaska, with the
Alaska breeding population currently listed as Threatened under the
Endangered Species Act. As such, frequent status assessments are needed
to determine correct listing status and to guide recovery actions.
Unfortunately, the rarity and low density of this eider in Alaska make
inferences of population size, trend, and distribution especially
difficult. Moreover, observations of Steller's eider on surveys are
extremely variable among years, with some years having no observed
eiders \citep{obritschkewitsch2008steller, wilson2025population}.
It has been hypothesized that these eiders forgo breeding in some years
when brown lemming (\emph{Lemmus trimucronatus}) abundance and
their avian predators are low \citep{quakenbush2004breeding}. Such
cycles or intermittent breeding would further complicate methods to
estimate long term trends or model population viability
\citep{dunham2017evaluating}. The effort described here is an attempt to
make the most of the available data using modern model-based methods to
distinguish sampling variation from the underlying signal of population
differences between years and locations.

The approach here is inspired by Miller et al \citep{miller2013} and
Amundson et al. \citep{amundson2019}, where generalized additive models
\citep{hastie1986generalized, wood2017} are used to estimate
temporal and spatial variation in eider density. This approach to
estimating population trend also borrows from the work of
\citep{sauer2002hierarchical, rosenberg2019decline, smith2021north} among others, although the focus here is only one species. In general, the advantage of model-based estimates are that they help to separate noise (i.e., sampling variation) from the true underlying signal (i.e., temporal or spatial patterns in density \citep{BDA3, hooten2015guide}, many others) and is closely related to the concept of shrinkage of random effects
\citep{efron1977stein, BDA3}, where information is shared across
samples in time, space, or other groups. This contrasts with pure
design-based estimates, which make minimal assumptions, but suffer from
high variability in this rare species. Specifically, sampling variation
alone might lead to observations of zero eiders even when eiders did
exist in the survey area. While correcting for the detection process is
possible for design-based estimates, in the case of zero observations,
no correction can be made, and these years will tend to under represent
the true animal density. Sampling variation also inflates the
year-to-year variance in estimates and will tend to make population
trend estimates less precise and potentially more extreme, especially
when multiple species are compared or statistical significance tests are
used to detect trends \citep{sauer2002hierarchical}. In contrast,
model-based estimates require the specification of a probability model
for observations and help distinguish sampling from non-sampling
variation. In the spatial and temporal context, where correlation is
generally high and positive between nearby measurements, models that
exploit this feature are especially helpful, and nearly all modern
methods to detect trends in time and space use such models (see
citations above, and \citep{kery2015applied, kery2020applied}).

Below, I apply spatio-temporal models to provide the best available
estimates of Steller's eider density and annual population size on the
Arctic Coastal Plain (hereafter, ACP) and a smaller but highly sampled
triangle-shaped area (Triangle survey) near Utqiagvik, Alaska. I first
describe the data (\S Survey areas and data source). I then describe the
calculations of design-based estimates (\S Design-based estimates). This
is the first report of design-based estimates for the Triangle data that
include an estimate of uncertainty. These provide a reference point
against model-based estimates. I then describe the data formatting and a
variety of spatio-temporal generalized additive models fit to the data
(\S Model-based estimates). Posterior simulations from the fitted model
estimates are then used to calculate total population size and trend
(\S Prediction and posterior simulation). I also describe the detection
estimate used during posterior simulations (\S Incorporating detection).
Finally (\S Wavelet analysis), I conduct a wavelet analysis of the
posterior mean population size across years to estimate and
statistically test for a periodic cycle in the eider population. This is
of practical significance because it has been hypothesized that nesting
effort of Steller's eiders is highest when brown lemming populations are
at or near peak of their boom-bust cycles, and if there is a strong
cycle, estimates of trend will depend on the relative endpoints used in
building the estimate. Thus, if the population is cyclic, consistent
relative positions in the cycle should be used for trend estimates.
Throughout the results I provide some interpretation and provide
comparisons to past work (\S Results and Interpretation). I then provide
some recommendations for model-based estimation of population size and
trends, compare these methods to related efforts, and discuss the
importance of the population cycles and detection estimates for species
status assessment under the US Endangered Species Act (\S Discussion,
\S Conclusion).

\section{Methods}\label{methods}

\subsection{Survey areas and data
source}\label{survey-areas-and-data-source}

Data used in this report came from the ACP Survey conducted by the
United States Fish and Wildlife Service, Division of Migratory Bird
Management, Alaska Region, and a survey conducted by ABR, Inc.-
Environmental Research and Services (ABR) near Utqiagvik, Alaska, under
contract by the United States Fish and Wildlife Service, Ecological
Services Field Office, Fairbanks, Alaska, (hereafter, the Triangle
survey, Fig~\ref{fig-area}). The ACP survey has been described in
\citep{amundson2019} and \citep{wilson2025population}. Unlike in
\citep{amundson2019}, here only data from the actual ACP survey (2007 to
2019 and 2022 to 2024) are used. Amundson et al. \citep{amundson2019} also included data from a different survey (before 2007) with different seasonal timing that causes a difference in observed response (a major subject of the effort in \citep{amundson2019}). Data collected before 2007 have also not undergone the same level of quality control, which was a major effort for this analysis. The ACP data from 2007 to 2024 used in this report, a description of quality control processes, and a description of data manipulations are available at \citep{osnas-github, osnasACP} (see \nameref{s1-acpdata}).

The Triangle survey is described in reports from ABR
\citep{obritschkewitsch2008steller, triangleABR}. Data for the
Triangle survey were obtained from ABR and are available at
\nameref{s2-tridata}. A lengthy quality control process was applied to
make data available for use in this report (documented in the file
\texttt{wrangle\_ABR.R}, available at \nameref{s3-code}).

Important differences between the surveys include the much higher
sampling coverage of the Triangle survey (25-50\%) compared to the ACP
survey (approximately 1 - 8\%), the much larger area of the ACP survey,
and the sampling effort is stratified in the ACP survey
(Fig~\ref{fig-area}). Hen eiders, birds occurring off of defined
transects, and behavior (flying/not flying) are also recorded on the
Triangle survey, but these are not generally recorded during the ACP
survey. In addition, the ACP survey records all waterbird species, but
the Triangle survey records a smaller subset (king {[}\emph{Somateria
spectabilis}{]} and spectacled {[}\emph{Somateria fischeri}{]} eiders in
addition to Steller's, other species in some years). Otherwise, the two
surveys follow standard aerial waterfowl survey protocols
\citep{cws1987}. No flocked Steller's eiders were observed on either
survey, and females outside of male-female pairs and all off-transect
observations were dropped from the Triangle data for all analyses. There
was one case of an "open 1" Steller's eider in the ACP data set (a
general category for a group of unknown sex or pair status that is not
appropriate for an observation of a single), and it was assumed to be a
single male.

\begin{figure}

\caption{\label{fig-area}\textbf{Arctic Coastal Plain (ACP) and Triangle survey areas.} \newline{East-west transect lines are shown in black for one year on the ACP. Strata with different transect sampling intensities for the ACP are shown by different fill colors. The Triangle area is shown as a stratum, but it is completely contained within the ACP High stratum. Transects for the Triangle survey are very close (800m in most years) so would appear as a solid color at this scale. On the Arctic Coastal Plain survey, transects are rotated north-south over a four year period to increase spatial coverage (not shown). All map data are products of the U.S. Government and are in the Public Domain.}}

\end{figure}%

All general data manipulation and analysis was done with program R
\citep{r-base} in RStudio \citep{RStudio2024} with the tidyverse
packages \citep{tidyverse}. This manuscript was produced using Quarto
(https://quarto.org/) and the PLOS One journal template found at
https://github.com/quarto-journals/plos in an effort to maximize
reproducibility. Code for this analysis, results, and production of this
manuscript is maintained at https://github.com/USFWS/STEI-estimates.

\subsection{Design-based estimates}\label{design-based-estimates}

I calculated a design-based estimate using formula R3 of Fewster et
al. \citep{fewster2009} (also see \citep{buckland2001}, p.~79)
modified for strip transects. A ratio estimator, which is more common
for strip-transect surveys \citep{cochran1977, williams2002, thompson2012}, was investigated but found to be poor (very high variance) because most observations of Steller's eiders are on short transects in the northern end of the triangular-shaped area and longer transects in the south rarely have observed eiders. The ratio estimator can provide lower variance estimates when there is a positive correlation between the two variables (here, animal count and transect length or area) \citep{cochran1977, thompson2012}. In this survey area, however, there is a negative correlation between count and transect length (or area). This makes the ratio estimator worse than a simple plot-based estimator or the R3 estimator of Fewster \citep{fewster2009}. The point estimate was calculated as

\begin{equation}\phantomsection\label{eq-1}{\hat N = A \frac{\Sigma_i n_i}{\Sigma_i a_i}}\end{equation}
and the variance of \(\hat N\) was calculated as
\begin{equation}\phantomsection\label{eq-2}{var(\hat N) = \left( \frac{A}{a} \right)^2 \frac{a}{k-1} \sum_i a_i\left( \frac{n_i}{a_i} - \frac{n}{a} \right)^2}\end{equation}
where \(A\) is the total area of the study area, \(a_i\) is the area of
strip transect \(i\), \(n_i\) is the number of encountered single males
or pairs on strip transect \(i\), and \(n\) is the total encounters on
the \(k\) surveyed strips, and \(a\) is the total area of the \(k\)
surveyed strips. It is arguable if a finite population correction,
\((1 - a/A)\), is appropriate, so I left it out as it is not used by
Fewster et al. \citep{fewster2009} or Buckland et al.
\citep{buckland2001} because observations on transects are not
repeatable due to movement of birds and the detection process. Fewster
et al. \citep{fewster2009} showed that Eq (2) can overestimate the
variance of systematic surveys when there is a strong gradient in
density, as there is in this area. In this specific case, however, the estimator above is much better than a standard ratio estimator.

\subsection{Model-based estimates}\label{model-based-estimates}

I used generalized additive models (GAMs, \citep{hastie1986generalized, wood2017, pedersen2019hierarchical}) to estimate eider density as a function of location and time. A GAM can be thought of as a generalized linear mixed model that fits smooth functions (splines) of covariates to predict the response, in addition to standard linear mixed model terms if specified. The optimal degree of smoothing is determined during model fitting through a model selection-like process. The smoothed temporal or spatial effect, thus, can account for autocorrelation in these domains \citep{wood2017, pedersen2019hierarchical, hefley2017basis}. GAMs are
widely used in animal density surface modeling (e.g., \citep{miller2013, amundson2019, hollmen2023climate} and many others).

To set up the data for model fitting, I divided sampled transects into 1
(Triangle survey) or 6 (ACP survey) km segments and assigned
observations of eiders (including observation of zero eiders) to segment
centroids. Segments on the boundary of the study area were often smaller
due to boundary issues. I then assumed a half-strip width of 200 m and
calculated an area for each segment. Total number of eider observations
(pairs and males) were summed for each segment and the coordinates (in Alaska Albers Equal Area projection, EPSG:3338) of the centroid were used for spatial location. Year was used as the covariate for temporal effects. The same procedure was used for both the entire ACP and for the Triangle survey, but for the ACP-only model (see below) the segment size for model fitting was 6 km. This was done for computational reasons (time) related to fitting many different models and to be consistent with \citep{amundson2019}. When Triangle and ACP data were combined into one model, a common 1 km segment size was used. All spatial data manipulation was done using the R package \texttt{sf} \citep{sf}.

I fit a variety of models to explore spatial and temporal effects. The
linear predictor for each model was:

\[M0: Count \sim s(X, Y)\] \[M1: Count \sim s(X, Y) + s(Year)\]
\[M2: Count \sim s(X, Y) + s(Year) + ti(X, Y, Year)\]
\[M3: Count \sim s(X, Y) + s(Year) + s_{re}(fYear)\]
\[M4: Count \sim s(X, Y) + s_{re}(fYear)\] In the above, \(Count\) is
the total number of pairs or males observed in a transect segment,
\(s()\) indicates a smooth function, \(ti()\) indicates a ``tensor
product smooth'' (a multidimensional smooth that allows the units of the dimensions to differ), and \(s_{re}()\) is a random effect in the usual
mixed model sense (in
\texttt{mgcv:\ s(...,\ bs\ =\ \textquotesingle{}re\textquotesingle{})}).
\(Year\) is calendar year as a continuous numeric variable, and
\(fYear\) is calendar year as a factor variable. Model \(M0\) is just a
smooth of location that does not change through time. Model \(M1\) is a
smooth of location and a separate smooth of year. This model means that
on the log scale the spatial smooth does not change its relative shape
through time but the overall height changes as a smooth function of
year. Model \(M2\) is a smooth of location, year, and a spatio-temporal
interaction between location and year. Model \(M2\) allows the shape of
the spatial pattern to change through time and allows the time trend to
depend on location. Model \(M3\) estimates a smooth of location, a
smooth of year, and a separate random effect of year, which allows
abrupt non-smooth deviations from an underlying smooth trend
\citep{smith2021north}. Model \(M4\) estimates a smooth of location and
treats year as a simple random effect. The model with year as a random
effect still imposes smoothing on the year effect, just as in the usual
mixed model, where year effects are pulled or ``shrunk'' toward an
overall grand mean. In the models above, each effect is written as an
``average'' or ``partial'' effect relative to the others and all contain
an intercept (not shown). Thus, \(s(X,Y) + s(Year)\) is an average
effect of location and an average effect of Year as a deviation from an
overall intercept. All models used a negative binomial response, a log
link function, and the log segment area as an offset, which controlled
for varying areas of segments near study area boundaries. Thus, the
model is estimating the expected response in 1 \(km^2\) of area. The
scale parameter of the negative binomial was estimated during model
fitting. Other response models were fit (Poisson, Tweedie) and all were
found to be substantially worse fitting than the negative binomial. I
used a `thin plate regression spline' (\texttt{s(...,\ bs\ =\ "tp")}
in \texttt{mgcv}, \citep{r-mgcv}) as the spline basis function for all
smooth effects in all models. All model fitting and prediction was done
using the R package \texttt{mgcv} \citep{r-mgcv}. Model diagnostics were
inspected using the residual simulation methods in R package
\texttt{DHARMa} \citep{DHARMa} and visualizations using the R package
\texttt{gratia} \citep{r-gratia}. Special attention was given to
quantile-quantile plots and over-dispersion and zero-inflation metrics.
Residual simulations suggested that the negative binomial distribution
appropriately modelled the large number of response zeros in these data.
To select the best model, the \texttt{AIC} function was used with the correct degrees of freedom for a GAM (see \texttt{?mgcv::AIC.gam} in \citep{r-mgcv}).

For modelling the ACP data alone, additional models were fit that
contained a random effect for observer. The response for these models
was observer-specific; therefore, each segment was only 200 m wide
(i.e., each side of the plane was a separate observation 200 m wide). An
additional five models were fit where each model above also contained an
observer effect, \(s_{re}(Observer)\). Observer effects were not
estimated for the Triangle data because they were not recorded in most
years. When I combined the Triangle data with the ACP data, I used a
fixed factor effect for survey (Triangle or ACP) to model any average
difference between surveys and used a model with spatial and temporal
smooths as in \(M1\). A model with a spatio-temporal interaction
(\(M2\)) was also explored for the combined Triangle and ACP data, but
rejected. The two data series do not overlap during the period of
1999-2006, so the interaction term was not well estimated and produced
widely variable population estimates during posterior simulations for
the years 1999 to 2006. During posterior simulation (see below), the
survey covariate was set to predict the Triangle survey because this is
the area where detection was estimated. R code giving the details of
model fitting can be found in the supplemental information
(\nameref{s3-code}).

\subsection{Prediction and posterior
simulation}\label{prediction-and-posterior-simulation}

I used the \texttt{predict.gam} function from \texttt{mgcv} package
\citep{r-mgcv, wood2017} to predict the expected density of
eiders across the study area (Triangle or ACP) in each year, from 1999
to 2023 for the Triangle and from 2007 to 2024 for the ACP. Note that
years with no data collection are included in these predictions. To do
this, the study area was gridded into 1 \(km^2\) or smaller
cells--smaller when the cells intersect the boundary of the study
area--and centroids and areas of each cell were calculated. A data set
was then created by replicating these point locations and areas for each
year. This large data set was then used in the \texttt{predict} function
along with the model object from the best fitting GAM model. For spatial
maps to represent relative differences in eider density, predictions
were made on the response scale (\texttt{predict} option
\texttt{type\ =\ "response"}) and the year effect was excluded
(\texttt{exclude\ =\ "s(Year)"} or
\texttt{exclude\ =\ c("s(Year)",\ "ti(X,Y,Year)")}). Predictions at cell
centroids were used for the entire cell, that is, the continuous smooth
density surface was rasterized into 1 \(km^2\) or smaller cells for
display in maps.

Results of fitting a GAM in \texttt{mgcv} provide a posterior
distribution of model parameters \citep{wood2017, miller2025bayesian}. Therefore, posterior simulation was used to calculate population totals over the study area for each year (see examples in help files for \texttt{mgcv::predict} in \citep{r-mgcv} or \citep{wood2017}, p.~342-343). Predictions were made once on the same grid and years as described above but \texttt{type\ =\ lpmatrix} was specified so that a design matrix, \(\textbf X\), was returned with one row for each prediction location-year and one column for each term in the model. Multivariate normal samples of the model parameters, \(\textbf b_i\) were then simulated using the fitted model parameter vector and variance-covariance matrix. For the Triangle study area, direct simulation from a multivariate normal distribution was used. For all simulations on the ACP study area, a Metropolis-Hastings algorithm was used to obtain samples because large areas of zero observations caused poor performance based on direct multivariate normal simulation using the parameter estimates and covariance matrix (see
\texttt{?mgcv::gam.mh} in \citep{r-mgcv} or \citep{gratiaPost}). A
posterior sample on the response scale was then calculated as
\begin{equation}\phantomsection\label{eq-3}{\textbf y_i = \textbf a exp(\textbf X \textbf b_i )}\end{equation}
where \(\textbf a\) is a vector of the area of each grid cell and
\(\textbf y_i\) is a vector of the expected responses. Note that
\(\textbf y_i\) is the expectation of the negative binomial distribution
and not an actual realization; thus, it is \(>0\) for all predictions.
The above simulation was repeated a large number of times (500) and
results were stored.

Various derived quantities of the posterior samples can also be
calculated. To find the expected population total in a given year, the
sampled vector can be summed over all cells for a given year. Let \(i\)
index the posterior sample, \(j\) index year, and \(k\) index the cell,
then a posterior sample for the expected population total in year \(j\)
is
\begin{equation}\phantomsection\label{eq-4}{\hat Y_{ij} = \sum_k y_{ijk}.}\end{equation}
Because the model was fit to pairs and single males, the above gives the
total ``indicated pairs''. To transform this to birds and ``indicated
birds'', the above posterior sample would be multiplied by 2. A
detection corrected population total for ``indicated birds'' can be
found as
\begin{equation}\phantomsection\label{eq-5}{\hat N_{ij} = 2 \hat Y _{ij}/d_{ij}}\end{equation}
if detection varies only by year, where \(d_{ij}\) is a posterior sample
of detection in year \(j\). If detection varies with location or other
covariates, then the adjustment needs to take place at a lower level of
\(y_{ijk}\).

The posterior trend on the log scale from year \(j\) to \(j+t\) can be
found as
\begin{equation}\phantomsection\label{eq-6}{T_{it} = log \left[ (\hat N_{ij+t}/\hat N_{ij})^{1/t} \right] = \frac{log(\hat N_{ij+t}) - log(\hat N_{ij})}{t}.}\end{equation}
Note that this measure of trend is identical to the slope parameter in a
``log-linear'' regression when the smooth of year in the GAM model,
\(s(Year)\), gives a straight line. The advantage of the GAM is that the
notion of trend can be generalized to non-linear cases where the trend
may be increasing, decreasing, or changing cyclically through time. The
posterior distribution for any quantity can then be displayed using a
histogram or summarized with various metrics. I summarized the posterior
with the mean and the 0.025, 0.5 (median), and 0.975 quantiles. For
trend, I only used detection corrected posterior estimates (see below).
Full detail of the posterior simulations can be found in the
supplemental information (\nameref{s3-code}).

\subsection{Incorporating detection}\label{incorporating-detection}

I used a preliminary detection rate estimated from the second year
(2018) of an on-going field study using artificial decoys to estimate
detection. The detection analysis was completed by Catherine Bradley
(USFWS, retired) and showed that detection rate in 2018 depended only on
distance from transect (see \nameref{s4-bradley}). Detection was also
estimated during 2017, the first year of the study, but field protocol
issues caused problems during this pilot year, and the protocol was
stabilized in 2018. Because distance was not available for observations
outside the decoy detection study area (a sub-set of the Triangle area),
I calculated the average `unconditional' detection rate over the
transect half-strip width (\nameref{s4-bradley}). In this
context, `unconditional' is the detection rate estimated from
double-observer sightability trials that include decoys not observed by
either observer (a `00' capture history). I also used this same
detection rate for the ACP because no better estimate for Steller's
eider is available. The detection rate used here should be viewed as
provisional until a more complete analysis can be conducted
incorporating data from across multiple years and crews. In any case, it
is meant to be the average detection rate over all observations,
including the covariates of distance, year, observer, sun angle, etc.,
and was in fact used for the Steller's eider Species Status Assessment
of 2025 \citep{steiSSA2025}.

To calculate an average detection rate, I assumed that eiders (detected
and undetected) were expected to be uniformly distributed with distance
from the transect, which is true given the design of the Triangle and
ACP surveys. I then averaged detection and the variance of the detection
estimate over each of the four distance bins in \nameref{s4-bradley}. Because distance bins were of equal width, detection was estimated as a factor of bin (i.e, a continuous distance function was not estimated), and no other covariates were found important, the mean and variance are a simple equally weighted average over the distance bins. In \nameref{s4-bradley} only the mean and ninety-five percent credible interval was given, so I approximated the standard error of the estimate in each bin by the range divided by 2*1.96, which is the typical number of standard deviations in the range of a symmetric credible interval.

This worked out to a detection rate of 0.307 with a standard deviation
of 0.092 or a Beta(7.41, 16.72). Note that this detection prior will
result in a large amount of uncertainty in the estimated number of
eiders. At detection rates of 0.16 and 0.47 (the fifth and ninety-fifth
percentile of the distribution, respectively), the eider population
estimate will be increased about 6- and 2-fold, respectively. Thus,
increased information on detection would be valuable for improving our
knowledge of the eider population size. Interestingly, this detection
rate prior has a very similar mean but is much more variable compared to
that used in the past based on expert judgment
\citep{dunham2017evaluating}.

The posterior population estimate for ``indicated breeding birds''
(single males and pairs), accounting for constant detection, was then
calculated as
\begin{equation}\phantomsection\label{eq-7}{\hat N_{ij} = 2\hat Y _{ij}/d_{i}.}\end{equation}
with \(d_i\) sampled from a Beta distribution with the mean and standard
deviation above. Thus, detection was assumed constant across years and
other covariates (e.g., crews, singles, pairs). If detection varies with
year, location or other covariates, then the adjustment needs to take
place at a lower level as in (Eq~\ref{eq-5}).

In general, quantities \(\hat Y\) are designated as ``indicated pairs''
because single males are assumed to ``indicate'' an observed female on
the nest (the male-female sex ratio is assumed equal). (Technically,
within the tradition of waterfowl biology, ``indicated pairs'' include
flocked drakes in groups of less than 5. However, none of these were
observed in either data set; thus, ``indicated breeding pairs'' and
``indicated pairs'' are the same.) When ``indicated pairs'' are
multiplied by 2, \(2\hat Y\), then the quantity becomes ``indicated
birds'' because each pair is 2 individuals. Hereafter, I use ``indicated
breeding birds'', and I use the word ``index'' to distinguish quantities
that are not corrected for detection. Thus, \(2\hat Y\) is the
``indicated breeding bird index,'' and \(2\hat Y/d\) is simply
``indicated breeding birds,'' to reflect that it is a population
estimate, at least conditional on the assumptions inherent with
``indicated.''

\subsection{Wavelet analysis}\label{wavelet-analysis}

In order to formally estimate the period of a repeating cycle in the
Steller's eider population, I used wavelet analysis to detect and
statistically test for the presence of a cycle. Wavelet analysis is
similar to a Fourier transform of a time series where the time domain is
expressed in the frequency domain. Unlike a basic Fourier transform,
however, a wavelet transform localizes the frequency information in time
so that changes in the frequency of a signal with time can be examined.
The main product of a wavelet analysis is a ``power spectrum'' (also
know as a spectrogram) that is a three dimensional surface showing the
``power'' (z-axis or a measure of the strength of a particular frequency
or period) in a time series as a function of the wave period (y-axis)
and time (x-axis). High ``power'' ridges across time correspond to a
large signal at that frequency (or period) in the time series. If the
period is changing through time, the position of the ridge will change
relative to the period axis as time increases. An accessible
introduction can be found in \citep{cazelles2008wavelet} and in the
documentation to the R package \texttt{WaveletComp} \citep{r-wavelet}. I
used the function \texttt{analyze.wavelet} from the \texttt{WaveletComp}
package \citep{r-wavelet} to calculate the wavelet properties of the
posterior mean population time series estimate from the model that
combined both Triangle and ACP data. I used the null hypothesis of white
noise and 200 bootstrap simulations to test for statistical departure
form white noise (a flat or uniform frequency spectrum) and to calculate
a confidence interval for the period. I then used the functions
\texttt{wt.image} to visualize the power spectrum and \texttt{wt.avg} to
visualize the average power spectrum across years. To find the full
posterior of the dominant wave frequency, I calculated the wavelet
properties of 500 posterior samples of the population time series from
the best fitting GAM model and saved the wave period at the maximum
average power across time for each sample. I then used the mean and 95\%
credible interval to summarize this posterior.

\section{Results and Interpretation}\label{results-and-interpretation}

\subsection{Observations of Steller's
eider}\label{observations-of-stellers-eider}

Observed locations of Steller's eiders across the study area are shown
in Fig~\ref{fig-observations}. Most observations are in the northern
coastal area of the Triangle and Teshekpuk area. No observations of
Steller's eiders have been made in the two southernmost ACP strata.
Observations from the Triangle survey abruptly stop at the southern edge
of the Triangle survey area, suggesting that eiders sometimes exist
south of that area more often than observations from either survey
suggest. Fig~\ref{fig-count} shows the count of Steller's eiders by year
for both the Triangle and ACP surveys.

\begin{figure}

\caption{\label{fig-observations}\textbf{Observed Steller's eider during the Arctic Coastal Plain (2007-2024) and Triangle (1999-2023) surveys.} \newline{Arctic Coastal Plain (ACP) strata are shown by black lines. Locations of eider observations by survey are given by the purple (ACP) or orange (Triangle) dots. All map data are products of the U.S. Government and are in the Public Domain.}}

\end{figure}%

\begin{figure}

\caption{\label{fig-count}\textbf{Number of Steller's eider observations by year for the Arctic Coastal Plain and Triangle surveys.}}

\end{figure}%

\subsection{Design-based estimates}\label{design-based-estimates-1}

Design-based estimates for the both survey areas are shown in
Fig~\ref{fig-design}. No survey was conducted in 2020 for the Triangle
survey or in 2020 and 2021 for the ACP. Note the higher mean and
variance of estimates for the ACP compared to the Triangle survey and
that estimates for the ACP often overlap zero. This is due to the larger
area and smaller sampling fraction of the ACP design. Design-based
estimates make minimal assumptions and provide an important check on
model-based estimates. Two important assumptions behind design-based
estimates are that (1) the survey design is unbiased (random or
systematic selection of sampling units, here transects) and (2) the
survey was implemented as designed. These assumptions are met in these
surveys. An additional often unstated assumption is that the response is
measured without error. In this case, measurement errors include a large
detection bias and, presumably less often, species misidentification.
Detection bias causes a lower mean response and increased variance in
response. Finally, design-based estimates are based on estimating a mean
response across all sampled transects, and then use statistical sampling
theory to derive estimator variance. This works well for common species
and large sample size (many transects), but for rare species that might
not be encountered on any transects during a sample, the estimate and
its variance are necessarily zero when there are no detection events,
even when the species might have existed in the survey area. This might
have been the case in 2009 on the Triangle area (Fig~\ref{fig-design}A) and in 2010 for all strata of the ACP (Fig~\ref{fig-design}B), and
was the case in 2, 10, 16 and 16 years out of 16 for the High,
Teshekpuk, Medium, and Low strata of the ACP, respectively.

\begin{figure}

\caption{\label{fig-design}\textbf{Design-based estimates for each survey.}\newline{(A) Triangle survey area and (B) Arctic Coastal Plain survey area. Point estimates for each year are given by black dots contected by black lines across years. Confidence intervals (95\%, orange band) have been truncated at zero. No detection correction was applied.}}

\end{figure}%

\subsection{Model-based Estimates:
Triangle}\label{model-based-estimates-triangle}

Models fit to data from the Triangle showed that the best model contains
a spatial smooth, a year smooth, and a spatio-temporal smooth (model M2,
Table~\ref{tbl-modelsTri}). The model with a separate space and time
smooth (M1) and all other models were worse (\(\Delta AIC >5\)). For all
results below, I used model M2.

\begin{table}

\caption{\label{tbl-modelsTri}\textbf{AIC table for models fit to Triangle data.}}

\centering{
\fontsize{12.0pt}{14.4pt}\selectfont
\begin{tabular*}{\linewidth}{@{\extracolsep{\fill}}llrrr}
\toprule
Model & Linear Predictor & df & AIC & $\Delta$AIC \\ 
\midrule\addlinespace[2.5pt]
M2 & s(X,Y) + s(Year) + ti(X,Y,Year) & 35.48 & 2876.13 & 0.00 \\ 
M1 & s(X,Y) + s(Year) & 33.63 & 2881.94 & 5.81 \\ 
M3 & s(X,Y) + s(Year) + s(fYear) & 39.89 & 2884.29 & 8.15 \\ 
M4 & s(X,Y) + s(fYear) & 38.18 & 2888.06 & 11.92 \\ 
M0 & s(X,Y) & 19.78 & 2979.11 & 102.98 \\ 
\bottomrule
\end{tabular*}
}
\begin{flushleft} Model structures are described in the main text.
\end{flushleft}
\end{table}%

The spatial, temporal, and spatio-temporal partial effects of model M2
are shown in Fig~\ref{fig-model}. The largest approximate p-value for
the smooth terms was 0.003 for the spatio-temporal smooth. Effective
degrees of freedom for the spatial (Fig~\ref{fig-model}A), temporal (Fig~\ref{fig-model}B), and spatio-temporal (Fig~\ref{fig-model}C) smooth
were 16.1, 12.1, and 2.1, respectively. Highest densities of Steller's
eiders were in the northern section of the Triangle and the lowest
densities were in the southeast, but as shown by the spatio-temporal
smooth (Fig~\ref{fig-model}C), density decreased through time in the
southeast and increased in the northwest of the Triangle. Note that the
spatio-temporal smooth has essentially been shrunken to a simple 2D
plane that is changing through time such that eider density is
increasing in the northwest and decreasing in the southeast. Full
\texttt{mgcv} model objects and parameter values can be found in the
supplemental information (\nameref{s5-results}). The temporal pattern
appears cyclic with a period of 5-7 years and no strong directional
trend (Fig~\ref{fig-model}B).

\begin{figure}

\caption{\label{fig-model}\textbf{Partial effect plots for the best model fit to the Triangle survey data.}\newline{Plots of spatial (A), temporal (B), and spatio-temporal interaction (C) partial effects from the best fitting model. In (B) the shaded region is 2 standard errors. In (A) and (C) black lines show contour lines of equal partial effect, and color shows increased (red) or decreased (blue) density relative to the intercept. In (C) the interaction is shown by 4 density surfaces at different time points from 1999 to 2023. Each panel is as in (A) but has been simplified to the bounding rectangle of the survey area. Note that as
time progresses, the upper left of each panel increases in density as
indicated by the shift from blue to red.}}

\end{figure}%

Spatial density of Steller's eider across the Triangle after removing
the effects of year and the space-by-year interaction (i.e., the average
or partial effect of location) is shown in Fig~\ref{fig-densityABR}.
Relatively high densities have occurred in the north and moderate
densities in the east and far southwest. With the spatio-temporal
effect, however, these areas of higher densities in the south have
decreased.

\begin{figure}

\caption{\label{fig-densityABR}\textbf{Predicted average density of Steller's eider in the Triangle survey.}\newline{Predicted density after the effects of year and space-by-year interaction have been removed.}}

\end{figure}%

Population estimates across the Triangle survey area by year for
indicated breeding birds without (Fig~\ref{fig-yearABR}A) and with (Fig~\ref{fig-yearABR}B) a detection correction are shown in Fig~\ref{fig-yearABR}. There is a substantial increase in the population estimates, the uncertainty, and in the skew of the posterior toward higher population estimates with the application of the a detection correction (Fig~\ref{fig-yearABR}B). This is due to the low mean detection rate and high uncertainty in the detection prior. Note that the model predicts the population in 2020 when no survey was conducted. The population appears to cycle on a period of 5-7 years (Fig~\ref{fig-yearABR}B); thus, causing the posterior n-year trend to fluctuate at the posterior mean (Fig~\ref{fig-trendABR}). Trend estimates are more precise for longer lags, but also fluctuate from positive to negative at the posterior mean (Fig~\ref{fig-trendABR}).

\begin{figure}

\caption{\label{fig-yearABR}\textbf{Posterior estimates of Steller's eider in the Triangle survey area, 1999-2023.}\newline{(A) without accounting for detection and (B) after applying a detection correction. The black line is the posterior mean, the dashed black line is the posterior median, and the orange band is the 95\% credible interval.}}

\end{figure}%

\begin{figure}

\caption{\label{fig-trendABR}\textbf{Posterior trend estimates for Steller's eider in the Triangle survey area, 1999-2023.} \newline{The y axis is the log of the geometric mean growth rate, and the x axis is the lag-year trend, the lag of 10 gives the 10-year trend from 2013 to 2023. The black line is the posterior mean, the dashed black line is the posterior median, and the orange band is the 95\% credible interval. The
horizontal thick black line is the y-axis origin.}}

\end{figure}%

\subsection{Model-based Estimates: ACP}\label{model-based-estimates-acp}

Models fit to the ACP data showed that a model with a random effect of
year, rather than a smooth of year, fit best
(Table~\ref{tbl-modelsACP}). In addition, models with a random observer
effect contributed nothing to improvement in AIC
(Table~\ref{tbl-modelsACP}). Upon examining model summary statistics,
models with a smooth of year shrunk to a straight line and observer
effect variance was essentially zero (results can be found in model
objects of \nameref{s5-results}). These results are likely due to the
few observations of Steller's eider on the ACP survey and the overall
sparse nature of the sampling. The year random effect of the best model
had the largest p-value (0.008) and an effective degrees of freedom of
7.9, indicating substantial shrinkage toward the mean of the 16 year
effects. The spatial smooth was substantially simpler (effective degrees
of freedom = 9.8) than the best model for the Triangle and showed a
simple decline in eider density as distance from the coast increased
(Fig~\ref{fig-modelACP}). Despite the shrinkage of the temporal random
effects, they still varied greatly (see below).

\begin{table}
\begin{adjustwidth}{-2.25in}{0in}
\caption{\label{tbl-modelsACP}\textbf{AIC table for model fit to Arctic Coastal Plain data.}}
\centering{
\fontsize{12.0pt}{14.4pt}\selectfont
\begin{tabular*}{\linewidth}{@{\extracolsep{\fill}}llrrr}
\toprule
Model & Linear Predictor & df & AIC & $\Delta$AIC \\ 
\midrule\addlinespace[2.5pt]
M4 & s(X,Y) + s(fYear) & 23.65 & 531.58 & 0.00 \\ 
M4.obs & s(X,Y) + s(fYear) + s(Observer) & 23.65 & 531.58 & 0.00 \\ 
M3 & s(X,Y) + s(Year) + s(fYear) & 24.51 & 531.83 & 0.25 \\ 
M3.obs & s(X,Y) + s(Year) + s(fYear) + s(Observer) & 24.53 & 531.86 & 0.28 \\ 
M0 & s(X,Y) & 14.54 & 537.44 & 5.86 \\ 
M0.obs & s(X,Y) + s(Observer) & 14.54 & 537.44 & 5.86 \\ 
M1 & s(X,Y) + s(Year) & 15.60 & 539.22 & 7.64 \\ 
M1.obs & s(X,Y) + s(Year) + s(Observer) & 15.60 & 539.22 & 7.64 \\ 
M2 & s(X,Y) + s(Year) +ti(X,Y,Year) & 14.07 & 540.47 & 8.89 \\ 
M2.obs & s(X,Y) + s(Year) + ti(X,Y,Year) + s(Observer) & 14.07 & 540.47 & 8.89 \\ 
\bottomrule
\end{tabular*}
}
\begin{flushleft}  Model structures are described in the main text. The suffix
`.obs' indicates a model that contained a random effect of observer.
\end{flushleft}
\end{adjustwidth}
\end{table}%

\begin{figure}

\caption{\label{fig-modelACP}\textbf{Partial effect of the spatial smooth for the Arctic Coastal Plain model M4.} \newline{The map of density shows decreasing density further to the south or away from the coast. Black lines show contour lines of equal partial effect, and color shows increased (red) or decreased (blue) density relative to the intercept.}}

\end{figure}%

Posterior simulations of the year-specific total across the whole ACP
are shown in Fig~\ref{fig-yearACP} without a detection correction (Fig~\ref{fig-yearACP}A) and with a detection correction (Fig~\ref{fig-yearACP}B). The population appears to fluctuate in a similar manner as in the model fit to Triangle survey data, but the year-specific estimates for the ACP are much less precise. Because the best model contained only a simple random effect to model year and the survey was not completed in 2020 and 2021, predictions were not made for these years using this model. Such predictions could be made by simulating random effects for these years but this would simply give an estimate that spans the range of historical estimates because there is no temporal correlation in this model.

\begin{figure}

\caption{\label{fig-yearACP}\textbf{Posterior estimates of Steller's eider across the Arctic Coastal Plain survey area, 2007-2024.} \newline{(A) without accounting for detection and (B) after applying a detection correction. The black line is the posterior mean, the dashed black line is the posterior median, and the orange band is the 95\% credible interval.}}

\end{figure}%

Note that year-specific estimates of the total number of eiders across
the ACP based on model \(M4\) (Fig~\ref{fig-yearACP}) are less extreme than those based on design-based estimates (Fig~\ref{fig-design}B). That is, the model-based estimates are lower in high abundance years and higher
in low abundance years compared to the design-based estimates. This is
due to the `shrinkage' or regularization of the year random effect. In
general, design-based estimates might over-estimate the eider population
when one or more eiders happen to be observed on a transect in one year
and under-estimate when one or more eiders are not observed on a
transect, even though they are in the study area.

Trends for the Steller's eider population across the ACP are shown in
Fig~\ref{fig-trendACP}. For all time ranges (lags), the trend was not
well estimated. For example, the 10-year trend posterior mean was 0.11
(95\% CI: -0.03, 0.23), which could be a modest decrease to a very fast
increase (cf., Fig~\ref{fig-yearACP}B). For all time periods, the 95\%
credible interval overlaps 0 and the upper bound was often
\textgreater{} 0.15 (Fig~\ref{fig-trendACP}).

\begin{figure}

\caption{\label{fig-trendACP}\textbf{Posterior trend estimates for Steller's eider in the Arctic Coastal Plain survey area, 2007-2024.} \newline{The y axis is the log of the geometric mean growth rate, and the x axis is the lag-year trend, the lag of 10 gives the 10-year trend from 2014 to 2024. The black line is the posterior mean, the dashed black line is the posterior median, and the orange band is the 95\% credible interval. The horizontal thick black line is the y-axis origin.}}

\end{figure}%

\subsection{Model-based Estimates: Triangle and ACP
combined}\label{model-based-estimates-triangle-and-acp-combined}

Partial effects of the model fit to combined data are shown in
Fig~\ref{fig-mapComb}. This model contained a spatial smooth, a temporal
smooth, and an effect of survey crew, but no spatio-temporal
interaction. All p-values for the smooth terms were very small
(< 0.001) and the survey crew effect p-value was
0.05. Effective degrees of freedom for the spatial and temporal smooths
were 19.2 and 12.8, respectively, indicating less shrinkage than the
ACP-only model and similar to the Triangle model. Similar to the
ACP-only model, eider density decreased from north to south and with
distance from the coast (Fig~\ref{fig-mapComb}A). Similar to the
Triangle model, the temporal trend fluctuated with a 5-7 year period
(Fig~\ref{fig-mapComb}B). The survey effect was estimated at 0.46 (SE =
0.23) on the log scale. Because the Triangle survey was defined as the
baseline (intercept), this means that the ACP crew observed
approximately sixty percent higher density at the posterior mean on
average as compared to the Triangle crew after controlling for year and
spatial location (95\% CI: 1.7\% - 147\%). Spatial predictions on the
response scale for 2013 are shown for the Triangle area in
Fig~\ref{fig-mapComb}C. Because there is no space-time interaction in
the model, the relative density surface is a simple multiplicative
relationship between all years as shown by the temporal partial effect
(Fig~\ref{fig-mapComb}B).

\begin{figure}

\caption{\label{fig-mapComb}\textbf{Partial effect plots for the best model fit to the combined Triangle and Arctic Coastal Plain data.} \newline{Partial effect plots for (A) spatial and (B) temporal effects. There is no space-time interaction in this model. (C) spatial prediction of eider density in 2013 zoomed in to the Triangle area with no
detection correction. Because there is no space-time interaction,
relative densities are the same across years. Similarly, the detection
correction does not vary across years or other covariates, so the relative
surface does not change with or without a detection correction.}}

\end{figure}%

Posterior distributions for the year-specific population predictions
over the entire ACP area are shown in Fig~\ref{fig-yearComb}A without a
detection correction and in Fig~\ref{fig-yearComb}B after the detection
correction is applied. Note that the predictions are made through
1999-2006 and 2021 when ACP survey data does not exist and for 2020 when
no data exists. This is possible because of the space-time structure of
the model. For 2020, predictions are due to the temporal correlation
(smooth) and interpolating from nearby years when data exists. In
1999-2006 and 2021, the existing Triangle data is used to inform the
temporal response, and then the temporally-constant spatial effect
informs the prediction across space.

\begin{figure}

\caption{\label{fig-yearComb}\textbf{Aunnual estimates of Steller's eider across the whole Arctic Coastal Plain, 1999-2024.} \newline{Estimates are from the best model fit to the Arctic Coastal Plain and Triangle data combined. Posteriors are given (A) without accounting for detection and (B) after applying a detection correction. The black line is the posterior mean, the dashed black line is the posterior median, and the orange band is the 95\% credible interval.}}

\end{figure}%

Posterior estimates of population trend show a strong increase over the
most recent years and fluctuating trends as the time period increases
(Fig~\ref{fig-trendComb}). The posterior estimate for the 10-year trend
is -0.03 (-0.15, 0.06) and for the entire series the 25-year trend is
-0.02 (-0.07, 0.02).

\begin{figure}

\caption{\label{fig-trendComb}\textbf{Posterior trend estimates for Steller's
eider on the Arctic Coastal Plain, 1999-2023.} \newline{Estimates are from the best model fit using data from both the Triangle and Arctic Coastal Plain surveys. The y axis is the log of the geometric mean growth
rate, and the x axis is the lag-year trend, the lag of 10 gives the 10-year trend from 2014 to 2024. The black line is the posterior mean,
the dashed black line is the posterior median, and the orange band is
the 95\% credible interval. The horizontal thick black line is the
y-axis origin.}}

\end{figure}%

\subsection{Wavelet analysis}\label{wavelet-analysis-1}

A wavelet power spectrum clearly revealed a high power ridge across
time that appears stationary with a period around 6 years
(Fig~\ref{fig-wavelet}). Another ridge of much lower power appeared at a
15 year period. When the spectrogram power is averaged across all years
the maximum power was found at period 6.53 (95\% CI: 5.46, 7.93). The
lower power ridge gave a maximum average power at 15.26 years with a
simulation pvalue of p = 0.07, suggesting weak evidence for a
longer period oscillation. When 500 posterior samples of the time series
were sampled from the GAM and a spectrogram found for each sample, the
posterior mean maximum power was found at period 6.52 (95\% CI: 6.23, 6.76), suggesting very high confidence in maximum average power at this period across the range of posterior samples of the GAM.

\begin{figure}

\caption{\label{fig-wavelet}\textbf{Wavelet power spectrum for Steller's eider across the Arctic Coastal Plain.} \newline{The power spectum was calculated at the posterior mean of the total population estimate.
Black points and lines give ridges where wavelet power is a local
maximum. White lines give a 95\% confidence interval for the period.
Larger power levels (red) correspond to greater probability density of
the wave period. The `cone of influence' is shown by the slightly
blurred area where boundary effects on estimating wave period are
large.}}

\end{figure}%

Such periodic cycles will cause any estimate of trend to depend on the
phase of the time points used for the calculation. Thus, when
calculating trend, it might be useful to control for period so that
start and end points reference the same relative position of the cycle.
That is, when using the results presented here for trend, use
approximate integer multiples of the period (e.g., 6.5, 13.0, 19.5, and
so on) so that similar phases of the cycle are compared. For example, a
19-year trend calculated from 2003 to 2022 (approximately trough to trough) gives a posterior mean of 0.03 (95\% CI: -0.01, 0.09).

\section{Discussion}\label{discussion}

Steller's eider numbers were estimated using generalized additive models
and traditional design-based methods for two survey areas, and then
using models after combining the two data sets. The Triangle data survey
is a small subset of the ACP that contains most of the eiders and is
sampled intensively; whereas, the ACP is a much larger area that covers
additional areas where Steller's eider have been observed, such as the
Teshekpuk area, but is sampled at a lower intensity. Because of the
rarity of Steller's eider in all areas but the most northern area of the
Triangle, design-based estimates are imprecise and the sampling process
can cause increased variance and zero observations in some years. This
is especially the case for the ACP survey where few or no Steller's
eider exist in most sample strata. Model-based estimates can ameliorate
some of these issues by using an explicit probability model for the
sampling process and sharing information across years and space to
smooth estimates. Such smoothing helps to remove sampling variation and
detect the spatio-temporal population dynamics, such as the periodic
cycles (Fig~\ref{fig-yearComb}), average spatial variation in density
across time (Fig~\ref{fig-densityABR}), or dynamics in spatial density
across time (Fig~\ref{fig-model}C).

Which results should one use? This depends on the purpose. If one is
simply interested in comparing survey results across years, then the
design-based estimate can help make rough assessments of changes in the
underlying population by comparing estimates and confidence intervals
across years. These calculations can also be useful for assessing design
properties of the surveys, such as sampling effort and power and are
very simple to do. While detection corrections can be applied to
design-based estimates, interpreting years of zero observations is
problematic, as no detection correction can be made in these years
without specifying a probability model that explicitly allows for zero
observation when animal are present. If one is interested in population
size, population trends, or variation in density across space, then
some form of probability model that includes ``regularization'' of
estimates should be used. While other models can be used, the GAMs
presented here are easily applied, seem to describe the data well, and
have all of the advantages discussed above. They are also commonly
applied for trend estimation (e.g. \citep{amundson2019, rosenberg2019decline, smith2021north}). Another advantage of the linear models used here (GAMs) is that interpretation and understanding is straightforward, where effects can be partitioned out into time, space, and space-by-time interaction effects (e.g.,
Fig~\ref{fig-model}).

For inference in the Triangle area, results from the models fit to the
Triangle-only data are best. This is for several reasons. First,
sampling intensity is high and consistent across space. This affects the
amount of smoothing that the GAM will select. Under sparse sampling, a
smoother surface will result, all else being equal (i.e., the true
wiggliness of the density surface). Compared to results in the Triangle
area from the combined model, the smooth of the density surface will
reflect the average properties across the whole ACP and this will result
in over-smoothing the Triangle area. This can be seen by comparing
Fig~\ref{fig-densityABR} to Fig~\ref{fig-mapComb}C. In addition, the
high sampling intensity across a relatively long time period has allowed
the estimation of a simple space-time interaction effect in the Triangle
area where density is increasing in the north and decreasing in the
south through time (Fig~\ref{fig-model}C). The model also estimates the
temporal effect well with much better precision than the design-based
estimates (compare Fig~\ref{fig-model}B to Fig~\ref{fig-design}A).

For inference outside the Triangle or the ACP as a whole, the model
using the combined data sets is best because design-based estimates
using only the ACP data are extremely imprecise for the Steller's eider
(Fig~\ref{fig-design}B) and even model-based estimates, while much
better, are not precise (Fig~\ref{fig-yearACP}) compared to results from
the combined data. Thus, combining data sets allows the Triangle data to
inform the temporal estimate, while the ACP data allows estimation of
density outside the Triangle area. Comparing Fig~\ref{fig-yearABR} to
Fig~\ref{fig-yearComb} shows that nearly half the population might be found outside the Triangle area in some years. Because the
space-time interaction of the Triangle-only model shows a decreasing
trend in the south (Fig~\ref{fig-model}C), there may be fewer eiders
outside the Triangle in recent years than there were in earlier years.
In any case, including the ACP survey data allows estimation of eider
density in the Teshekpuk and other areas of the ACP. If the ACP data is
not available, however, then survey effort in the Triangle area only is
likely sufficient to monitor abundance and trends of the population,
given the current distribution.

The model and approach used here is similar to that of Amundson et al. 
\citep{amundson2019}, with two important simplifications. First, the
time series for the ACP data was restricted to 2007 through 2024. This
allows a much simpler model because adjusting for survey timing
differences pre-2007 is not necessary. Thus, the phenology covariates
used in \citep{amundson2019} are not included or necessary in the models
used here. Second, restricted maximum likelihood as implemented in
\texttt{mgcv} \citep{wood2011fast} was used instead of direct Bayesian
simulations (Markov chain Monte Carlo). This is much easier from a
development and computational time perspective, making the current
approach easy to apply without custom coding in Bayesian software.
Instead, these models can be applied in R using more familiar linear
model syntax, the model fitting algorithm is much faster than direct
Bayesian simulations, and results are easier to assess. Another
difference is that the current effort estimated observer effects, which
were not in \citep{amundson2019}, and incorporated a detection estimate.
While the observer effects contributed nothing for the Steller's eider
(Table~\ref{tbl-modelsACP}), they have for other species of waterbirds
in separate analyses \citep{wilson2025population} or other systems
\citep{sauer1994observer, fink2023double}. Therefore, observer
effects are probably also important for the Steller's eider but simply
cannot be justified based on parsimony inherent in AIC model selection
criteria. The effect of ``Survey'' in the combined analysis is curious.
This suggests a strong effect of observer if there are not large
differences in protocol, which there are not, or in implementation
between the surveys. This effect deserves more investigation. In any
case, for population estimates, predictions were made setting the survey
parameter to the Triangle (the intercept) because this is the survey
area for which detection was estimated. Because there are so few
observations of Steller's eider on the ACP survey, however, this effect
may not be estimated well and deserves more investigation. One solution
might be to sample the Triangle area with higher intensity during the
ACP survey.

One interesting comparison to Amundson et al. \citep{amundson2019} is that they found very linear time trends for Steller's eider. This was also the case here when models included a spatio-temporal interaction and were fit to ACP-only data (model M2, Table~\ref{tbl-modelsACP}, results in
\nameref{s5-results}). This model, however, was not well-supported
compared to simpler models (Table~\ref{tbl-modelsACP}). When fit to data
that combined ACP and Triangle data, the model with a spatio-temporal
interaction was well-supported but was not used here because it gave
widely unbelievable predictions from 1999 to 2006. I attributed this to
the spatio-temporal mismatch in data sets during this period and a
spatio-temporal mismatch in sampling intensity. This was not a problem
in the data of \citep{amundson2019}, but they did not fit simpler models
without a space-time interaction. Because the simpler model without a
space-time interaction reported here produced reasonable predictions and
all other data sets produced similar cyclic temporal patterns, I believe
the simpler model with temporal and spatial effects is a more
appropriate model when data sets are combined, and the results of
\citep{amundson2019} are simply due to a lack of power to detect
anything other than linear smooths of year in their more complex model.
In any case, a perfectly linear (on the log scale) temporal trend is
hard to believe and most likely represents a modelling artifact of this
more complex model when confronted with sparse or mismatched data in
time and space.

\section{Conclusion}\label{conclusion}

Steller's eider populations on the ACP, including the Triangle area, are
clearly undergoing cyclical patterns of abundance with a period of about
6.5 years (Fig~\ref{fig-yearABR}, Fig~\ref{fig-yearComb},
Fig~\ref{fig-wavelet}). While it has been hypothesized that this is due
to lemming cycles \citep{quakenbush2004breeding}, this effort cannot
determine the cause of the cycles. Lemmings tend to cycle with a shorter
period, closer to 3-6 years with a median period around the arctic of
3.7 years \citep{gauthier2024taking}. The cycle near the Triangle survey
is somewhat unclear from historical data \citep{pitelka2018demography},
so unless the lemming period is \(\ge\) 5.5 years, the estimate of
periodicity presented here does not support a simple concordance between
eider and lemmings or their predators. It is clear that eider cycles
exist and that the population has been fluctuating around stable cycles
for at least 25 years (Fig~\ref{fig-wavelet}). Thus, any measure of trend will show an increasing, decreasing, or stable estimate depending on the beginning and ending points for the trend calculation (Fig~\ref{fig-trendABR}, Fig~\ref{fig-trendComb}). Recovery criteria and species status
assessments should incorporate this understanding into trend
determinations. Attention should be given to the start and end points
for trends so that similar points in the cycle are used (integer multiples of approximately 6.5). While each cycle is a stochastic realization that
will not follow exactly the same period, care should be taken when
interpreting estimates of trend. For example, a 10-year trend from 2013
to 2023 would show a large and significant decrease simply due to the
cycle phase; whereas, 2009 to 2019 would show a large increase. Instead,
the population seems to be undergoing fairly stable cycles of constant
period.

The detection estimate used here (0.307, SD = 0.092) is substantially
lower and more variable than used in the Species Status Assessment of
2019 (0.43, SD = 0.028, \citep{steiSSA2019}), which was based on
estimates for long-tailed duck (\emph{Clangula hyemalis}) and did not account for non-detection by both observers (a `00' capture history). Analyses previous to 2019 (USFWS, unpublished) used a point estimate of 0.3 (SD = 0) based on expert judgement, which is remarkably close to the mean estimate used here, but did not include any uncertainty. Dunham and Grand \citep{dunham2017evaluating} used the same point estimate and allowed for some uncertainty (SD = 0.02). Therefore, the posterior mean
population estimates reported here will be higher and skewed toward
larger values than past estimates. Better estimates of detection could
have large effects on the population estimates. Until such estimates are
available and population estimates are updated, the current estimates
should be treated as the best available, and the approach used here
seems a reasonable way to incorporate prior (or posterior) distributions
of detection into population estimates.

\section{Acknowledgements}\label{acknowledgements}

Comments from Annie Maliguine, John Nash, Dave Safine, Julian Fischer,
Amy Pocewicz, and Chuck Frost improved the manuscript. Cat Bradley
designed and analyized the initial years of the detection experiment.
Tim Obritschkewitsch (ABR, Inc.) provided data, including improvements
in data over that available in 2019, and helped to interpret data so
that it could be used in this analysis. The findings and conclusions in
this article are those of the author and do not necessarily represent
the views of the U.S. Fish and Wildlife Service.

\begin{thebibliography}{10}

\bibitem{obritschkewitsch2008steller}
Obritschkewitsch T, Ritchie RJ, King JG.
\newblock Steller's Eider Survey Near Barrow, Alaska, 2008.
\newblock ABR, Incorporated--Environmental Research and Services; 2008.
\newblock Available from: \url{https://catalog.epscor.alaska.edu/dataset/77d0ac11-90a2-4211-9183-237ec41b919d/resource/8ed2c5f0-1997-4701-983f-d9c221b97741/download/2008-barrow-stellers-eider-survey_final.pdf}.

\bibitem{wilson2025population}
Wilson HM, Safine DE, Frost CJ, Osnas EE.
\newblock Population indices, trends, and distribution of breeding waterbirds on the arctic coastal plain, Alaska, 2007-2024. 2025.
\newblock \href {http://dx.doi.org/https://doi.org/10.7944/dqf4-2z27} {doi:https://doi.org/10.7944/dqf4-2z27}.

\bibitem{quakenbush2004breeding}
Quakenbush L, Suydam R, Obritschkewitsch T, Deering M.
\newblock Breeding biology of Steller's eiders (Polysticta stelleri) near Barrow, Alaska, 1991-99.
\newblock Arctic. 2004:166-82.

\bibitem{dunham2017evaluating}
Dunham K, Grand JB.
\newblock Evaluating models of population process in a threatened population of Steller's eiders: a retrospective approach.
\newblock Ecosphere. 2017;8(3):e01720.

\bibitem{miller2013}
Miller DL, Burt ML, Rexstad EA, Thomas L.
\newblock Spatial models for distance sampling data: recent developments and future directions.
\newblock Methods in Ecology and Evolution. 2013;4(11):1001-10.

\bibitem{amundson2019}
Amundson CL, Flint PL, Stehn RA, Platte RM, Wilson HM, Larned WW, et~al.
\newblock Spatio-temporal population change of Arctic-breeding waterbirds on the Arctic Coastal Plain of Alaska.
\newblock Avian Conservation \& Ecology. 2019;14(1).

\bibitem{hastie1986generalized}
Hastie T, Tibshirani R.
\newblock Generalized additive models.
\newblock Statistical science. 1986;1(3):297-310.

\bibitem{wood2017}
Wood SN.
\newblock Generalized Additive Models: An Introduction with R, Second Edition.
\newblock Boca Raton, Florida: Chapman and Hall/CRC; 2017.

\bibitem{sauer2002hierarchical}
Sauer JR, Link WA.
\newblock Hierarchical modeling of population stability and species group attributes from survey data.
\newblock Ecology. 2002;83(6):1743-51.

\bibitem{rosenberg2019decline}
Rosenberg KV, Dokter AM, Blancher PJ, Sauer JR, Smith AC, Smith PA, et~al.
\newblock Decline of the North American avifauna.
\newblock Science. 2019;366(6461):120-4.

\bibitem{smith2021north}
Smith AC, Edwards BP.
\newblock North American Breeding Bird Survey status and trend estimates to inform a wide range of conservation needs, using a flexible Bayesian hierarchical generalized additive model.
\newblock The Condor. 2021;123(1):duaa065.

\bibitem{BDA3}
Gelman A, Carlin JB, Stern HS, Dunson DB, Vehtari A, Rubin DB.
\newblock Bayesian data analysis.
\newblock CRC press; 2014.

\bibitem{hooten2015guide}
Hooten MB, Hobbs NT.
\newblock A guide to Bayesian model selection for ecologists.
\newblock Ecological monographs. 2015;85(1):3-28.

\bibitem{efron1977stein}
Efron B, Morris C.
\newblock Stein's paradox in statistics.
\newblock Scientific American. 1977;236(5):119-27.

\bibitem{kery2015applied}
K{\'e}ry M, Royle JA.
\newblock Applied Hierarchical Modeling in Ecology: Analysis of distribution, abundance and species richness in R and BUGS: Volume 1: Prelude and Static Models.
\newblock Academic Press; 2015.

\bibitem{kery2020applied}
K{\'e}ry M, Royle JA.
\newblock Applied Hierarchical Modeling in Ecology: Analysis of distribution, abundance and species richness in R and BUGS: Volume 2: dynamic and advanced models.
\newblock Academic Press; 2020.

\bibitem{osnas-github}
Osnas EE. ACP-Mapping. GitHub; 2024.
\newblock \url{https://github.com/USFWS/ACP-Mapping}.

\bibitem{osnasACP}
Osnas EE. Arctic Coastal Plain Waterfowl and Waterbird Spatial and Temporal Trends.. U.S. Fish and Wildlife Service, Alaska Region, Migratory Bird Management Alaska; 2024.
\newblock \url{https://doi.org/10.7944/qtgn-y170}.

\bibitem{triangleABR}
Obritschkewitsch T, Bankert AR.
\newblock Steller's Eider surveys near Utqiaġvik, Alaska, 2023; 2024.

\bibitem{cws1987}
USFWS, CWS.
\newblock Standard Operating Procedures for Aerial Waterfowl Breeding and Ground Population and Habitat Surveys in North America; 1987.

\bibitem{r-base}
{R Core Team}. R: A Language and Environment for Statistical Computing. Vienna, Austria; 2024.
\newblock Available from: \url{https://www.R-project.org/}.

\bibitem{RStudio2024}
{Posit team}. RStudio: Integrated Development Environment for R; 2024.
\newblock Available from: \url{http://www.posit.co/}.

\bibitem{tidyverse}
Wickham H, Averick M, Bryan J, Chang W, McGowan LD, François R, et~al.
\newblock Welcome to the {tidyverse}.
\newblock Journal of Open Source Software. 2019;4(43):1686.
\newblock \href {http://dx.doi.org/10.21105/joss.01686} {doi:10.21105/joss.01686}.

\bibitem{fewster2009}
Fewster RM, Buckland ST, Burnham KP, Borchers DL, Jupp PE, Laake JL, et~al.
\newblock Estimating the encounter rate variance in distance sampling.
\newblock Biometrics. 2009;65(1):225-36.

\bibitem{buckland2001}
Buckland ST, Anderson DR, Burnham KP, Laake JL, Borchers DL, Thomas L, et~al.
\newblock Introduction to distance sampling: estimating abundance of biological populations.
\newblock Oxford (United Kingdom) Oxford Univ. Press; 2001.

\bibitem{cochran1977}
Cochran WG.
\newblock Sampling Techniques.
\newblock New York, New York: John Wiley and Sons; 1977.

\bibitem{williams2002}
Williams BK, Nichols JD, Conroy MJ.
\newblock Analysis and Management of Animal Populations.
\newblock New York, New York: Academic Press; 2002.

\bibitem{thompson2012}
Thompson SK.
\newblock Sampling.
\newblock Hoboken, New Jersey: John Wiley and Sons; 2012.

\bibitem{pedersen2019hierarchical}
Pedersen EJ, Miller DL, Simpson GL, Ross N.
\newblock Hierarchical generalized additive models in ecology: an introduction with mgcv.
\newblock PeerJ. 2019;7:e6876.

\bibitem{hefley2017basis}
Hefley TJ, Broms KM, Brost BM, Buderman FE, Kay SL, Scharf HR, et~al.
\newblock The basis function approach for modeling autocorrelation in ecological data.
\newblock Ecology. 2017;98(3):632-46.

\bibitem{hollmen2023climate}
Hollmen TE, Flint PL, Ulman SE, Wilson HM, Amundson CL, Osnas EE.
\newblock Climate change and coastal wetland salinization: Physiological and ecological consequences for Arctic waterfowl.
\newblock Functional Ecology. 2023;37(7):1884-96.

\bibitem{sf}
Pebesma E.
\newblock {Simple Features for R: Standardized Support for Spatial Vector Data}.
\newblock {The R Journal}. 2018;10(1):439-46.
\newblock Available from: \url{https://doi.org/10.32614/RJ-2018-009}. \href {http://dx.doi.org/10.32614/RJ-2018-009} {doi:10.32614/RJ-2018-009}.

\bibitem{r-mgcv}
Wood SN. mgcv: Mixed GAM Computation Vehicle with Automatic Smoothness Estimation; 2021.
\newblock R package version 1.8-36.
\newblock Available from: \url{https://CRAN.R-project.org/package=mgcv}.

\bibitem{DHARMa}
Hartig F. DHARMa: Residual Diagnostics for Hierarchical (Multi-Level / Mixed) Regression Models; 2022.
\newblock R package version 0.4.6.
\newblock Available from: \url{https://CRAN.R-project.org/package=DHARMa}.

\bibitem{r-gratia}
Simpson GL. {gratia}: Graceful {ggplot}-Based Graphics and Other Functions for {GAM}s Fitted using {mgcv}; 2023.
\newblock R package version 0.8.1.
\newblock Available from: \url{https://gavinsimpson.github.io/gratia/}.

\bibitem{miller2025bayesian}
Miller DL.
\newblock Bayesian views of generalized additive modelling.
\newblock Methods in Ecology and Evolution. 2025;16(3):446-55.

\bibitem{gratiaPost}
Simpson GL. {gratia}: Posterior simulation; 2024.
\newblock Version 0.9.2.9011.
\newblock Available from: \url{https://gavinsimpson.github.io/gratia/articles/posterior-simulation.html}.

\bibitem{steiSSA2025}
{U S  Fish and Wildlife Service}.
\newblock Status assessment of the {Alaska-breeding} population of {Steller’s Eiders}.
\newblock Fairbanks, {AK}: Fairbanks Fish and Wildlife Field Office; 2025. Version 2.
\newblock Available from: \url{https://iris.fws.gov/APPS/ServCat/DownloadFile/275532}.

\bibitem{cazelles2008wavelet}
Cazelles B, Chavez M, Berteaux D, M{\'e}nard F, Vik JO, Jenouvrier S, et~al.
\newblock Wavelet analysis of ecological time series.
\newblock Oecologia. 2008;156(2):287-304.

\bibitem{r-wavelet}
Roesch A, Schmidbauer H. {WaveletComp}: Computational Wavelet Analysis; 2018.
\newblock R package version 1.1.
\newblock Available from: \url{https://CRAN.R-project.org/package=WaveletComp}.

\bibitem{wood2011fast}
Wood SN.
\newblock Fast stable restricted maximum likelihood and marginal likelihood estimation of semiparametric generalized linear models.
\newblock Journal of the Royal Statistical Society Series B: Statistical Methodology. 2011;73(1):3-36.

\bibitem{sauer1994observer}
Sauer JR, Peterjohn BG, Link WA.
\newblock Observer differences in the North American breeding bird survey.
\newblock The Auk. 1994;111(1):50-62.

\bibitem{fink2023double}
Fink D, Johnston A, Strimas-Mackey M, Auer T, Hochachka WM, Ligocki S, et~al.
\newblock A double machine learning trend model for citizen science data.
\newblock Methods in Ecology and Evolution. 2023;14(9):2435-48.

\bibitem{gauthier2024taking}
Gauthier G, Ehrich D, Belke-Brea M, Domine F, Alisauskas R, Clark K, et~al.
\newblock Taking the beat of the Arctic: are lemming population cycles changing due to winter climate?
\newblock Proceedings of the Royal Society B. 2024;291(2016):20232361.

\bibitem{pitelka2018demography}
Pitelka FA, Batzli GO.
\newblock Demography and condition of brown lemmings (Lemmus trimucronatus) during cyclic density fluctuations near Barrow, Alaska.
\newblock In: Annales Zoologici Fennici. vol.~55. BioOne; 2018. p. 187-236.

\bibitem{steiSSA2019}
{U S  Fish and Wildlife Service}.
\newblock Status assessment of the {Alaska-breeding} population of {Steller’s Eiders}.
\newblock Fairbanks, {AK}: Fairbanks Fish and Wildlife Field Office; 2019. Version 1.
\newblock Available from: \url{https://iris.fws.gov/APPS/ServCat/DownloadFile/163633}.

\end{thebibliography}


\section{Supporting information}\label{supporting-information}

\paragraph*{S1 ACP Data}
\label{s1-acpdata}
{\textbf{ACP data used in these analyses.} The data can be found at
\url{https://doi.org/10.7944/qtgn-y170} \citep{osnasACP}.}

\paragraph*{S2 Triangle Data}
\label{s2-tridata}
{\textbf{Triangle data used for this report.} The data is maintained by
the author and is available at
\url{https://doi.org/10.7944/0x1v-4b77}.
The Triangle data used in this report was obtained from ABR, Inc.~}

\paragraph*{S3 R Code}
\label{s3-code}
{\textbf{R code to reproduce these analyses, results, and this report.}
All code can be found at the GitHub repository,
\url{https://github.com/USFWS/STEI-estimates}.}

\paragraph*{S4 Table 3}
\label{s4-bradley}
{\textbf{Detection estimates by distance from transect.} This is an unpublished USFWS report that describes the field experiment and analysis to estimate detection probabilty of Steller's eider by an aerial observer in 2017 and 2018. Table 3 is used to calculate the detection probability averaged over distance.}

\paragraph*{S5 Results}
\label{s5-results}
{\textbf{R model objects, posteriors samples, figures, and csv files
summarizing various posteriors.} All objects can be found at
\url{https://doi.org/10.7944/22x5-t856}.}


\nolinenumbers


\end{document}
